\subsection{Bit layout and stability for the Walsh--Hadamard transform}
\label{sec:binary}

At radix two, each recursive call transforms selected bits of an
address and leaves the other bits available for independent inputs.
The tape multiplication algorithm in \cite{oaT} uses this layout.
The same template also computes the normalized Walsh--Hadamard
transform. Its normalization can be absorbed into the recursive
calls, giving a fixed coefficient set and a bound on the accumulation
of numerical error.

\begin{proposition}[Binary implementations]\label{prop:binary-template}
Fix \(\theta>1/2\) and the arity \(A\), batch size \(W\) and call count
\(\sigma\) of Equation~\eqref{eq:kernel-template-saving} at \(\ell=2\).
\begin{enumerate}
\item An unnormalized batch on \(Af\) bits uses, outside its recursive
calls, pointwise linear maps on a fixed number of row entries, row
rewirings, and invertible \(\mathbb F_2\)-linear reindexings of the
\(A\) bits within each column. Its coefficients are drawn from a fixed
finite list of dyadic rationals and from signed powers of two with
exponent of magnitude \(O(f+1)\). Each call transforms a single
\(f\)-bit axis independently at every setting of the other bits.

The \(A\) axes may occupy contiguous slots consisting of \(f\) blocks
of \(K\) bits each. For any fixed \(0\le\rho<K\), the template can use
just bit \(\rho\) of each block. All other bits retain their addresses.
\item For even \(k\), the normalized transform \(2^{-k/2}D_2^{\otimes k}\)
has a uniform dyadic linear circuit of size
\(O_\theta(2^k(1+k)^\theta)\) with coefficients in a fixed finite
list outside the recursive calls. There are constants \(B,C,c>0\)
depending only on the template with the following property. If each
gate value is multiplied by \(1+\delta_g\), where
\(\lvert\delta_g\rvert\le u\le c(1+k)^{-B}\), the perturbed and exact
maps differ in Euclidean operator norm by at most \(Cu(1+k)^B\).
\end{enumerate}
\end{proposition}

\begin{proof}
The coefficient ring \(R_2=\mathbb Z[1/2]\) makes the fixed row
operations dyadic after descent. Table lookups use sums modulo two
of fixed column functions, and the phases are signs. The remaining
coefficients are powers of two from the axis normalizations and
from recovering the unnormalized target from the product \(F\).

By binary elimination each column reindexing is a fixed sequence of
additions from one axis to another, with an exchange implemented by
three additions. If \(x\) and \(y\) are the source and destination
slots for one such addition, the update is
\begin{equation}\label{eq:reindex}
 y_{\rho+iK}\longleftarrow y_{\rho+iK}+x_{\rho+iK}
       \quad\text{in }\mathbb F_2,\qquad 0\le i<f .
\end{equation}
Only the selected bits change. Fixing \(s\) inactive bits divides
one array of \(Z\) entries into \(2^s\) independent streams of
\(Z/2^s\) entries each. Group these streams into batches of \(W\),
padding the last batch. The recursive calls process in total
\[
 \sigma\lceil2^s/W\rceil\, Z/2^s
       \ \le\ (\sigma/W)(1+W/2^s) Z
\]
entries. This bound continues to hold for every slot width and offset.

For the normalized transform, take \(f\) to be the largest even
integer at most \((k-r)/A\) whenever this yields a positive child
width. The remaining \(k-Af<r+2A\) bits have even width and can be
handled by normalized four-point butterflies, as can the bounded
base cases. The padding bound and the strict saving in the recurrence
still apply.

Use a normalized transform in each axis call. For a common kernel
the \(2D\) calls supply the entire factor
\((2^{-f/2})^{2D}=2^{-Df}\). For a free rectangle the two partial
transforms together use those same \(2D\) axes. In a full spectral
conversion the common normalization commutes through its row table
and cancels with that of the inverse. Finally, \(F\) is already the
normalized target up to phases and a permutation. Thus all scalars
outside the recursive calls and the bounded remainder belong to a
fixed finite dyadic list.

Give the concatenated work rows their Euclidean norm. Each pointwise
stage is a direct sum of fixed-size circuits with bounded coefficients,
so its exact operator norm is bounded and perturbing its gates changes
the stage by \(O(u)\) for \(u\le1\). Permutations, zero padding and
restriction have norm at most one. Each recursive stage acts on
independent fibers by unitary transforms. Its perturbation has norm
bounded by the largest error on a fiber, even if different gates have
different errors. The number of successive stages per recursion level
is fixed. Telescoping their product therefore bounds the worst error
\(e(k)\) by \(C_0(u+e(f))\) whenever \(e(f)\le1\), with a constant
that includes the bounded remainder and base cases. The recursion
has depth \(O(1+\log_A(1+k))\). Iteration gives
\(e(k)\le Cu(1+k)^B\), and choosing \(c\) sufficiently small ensures
the child bound \(e(f)\le1\) at every level.
\end{proof}

Over \(\mathbb Q(\sqrt2)\) one can normalize odd widths as well.
Remove one bit from an odd-width call, apply \(2^{-1/2}D_2\) to that
bit, and use the even-width circuit on the rest. One may equally
normalize every axis call in a one-level template with arbitrary \(f\),
since its \(2D\) factors still multiply to \(2^{-Df}\). This retains
the call count, the inactive coordinates and the columnwise
reindexings. The scalars outside the calls now form a fixed finite
algebraic list. They can be approximated to \(p\) binary places in
time polynomial in \(p\) using dyadic arithmetic and the positive
root of \(2x^2=1\).
